\documentclass[10pt]{article}
\usepackage[utf8]{inputenc}
\usepackage[T1]{fontenc}
\usepackage{amsmath}
\usepackage{amsfonts}
\usepackage{amssymb}
\usepackage[version=4]{mhchem}
\usepackage{stmaryrd}
\usepackage{mathrsfs}
\usepackage{bbold}
\usepackage{graphicx}
\usepackage[export]{adjustbox}
\graphicspath{ {./images/} }
\usepackage{caption}
\usepackage{hyperref}
\hypersetup{colorlinks=true, linkcolor=blue, filecolor=magenta, urlcolor=cyan,}
\urlstyle{same}

\begin{document}
\captionsetup{singlelinecheck=false}
15.1. Hermite polynomials ..... 85\\
15.2. The Vandermonde determinant and the Hermite kernel ..... 88\\
15.3. Determinants of reproducing kernels ..... 90\\
15.4. The averaged empirical eigenvalue distribution ..... 92\\
16. Fluctuations of the Largest Eigenvalue of $G U E_{n}$, and Hypercontractivity ..... 94\\
16.1. $L^{p}$-norms of Hermite polynomials ..... 94\\
16.2. Hypercontractivity ..... 95\\
16.3. (Almost) optimal non-asymptotic bounds on the largest eigenvalue ..... 100\\
16.4. The Harer-Zagier recursion, and optimal tail bounds ..... 102\\
17. The Tracy-Widom Law and Level Spacing of Zeroes ..... 107

For the next ten weeks, we will be studying the eigenvalues of random matrices. A random matrix is simply a matrix (for now square) all of whose entries are random variables. That is: $X$ is an $n \times n$ matrix with $\left\{X_{i j}\right\}_{1 \leq i, j \leq n}$ random variables on some probability space $(\Omega, \mathscr{F}, \mathbb{P})$. Alternatively, one can think of $X$ as a random vector in $\mathbb{R}^{n^{2}}$ (or $\mathbb{C}^{n^{2}}$ ), although it is better to think of it as taking values in $M_{n}(\mathbb{R})$ or $M_{n}(\mathbb{C})$ (to keep the matrix structure in mind). For any fixed instance $\omega \in \Omega$, then, $X(\omega)$ is an $n \times n$ matrix and (maybe) has eigenvalues $\lambda_{1}(\omega), \ldots, \lambda_{n}(\omega)$. So the eigenvalues are also random variables. We seek to understand aspects of the distributions of the $\lambda_{i}$ from knowledge of the distribution of the $X_{i j}$. To begin, in order to guarantee that the matrix actually has eigenvalues (and they accurately capture the behavior of the linear transformation $X$ ), we will make the assumption that $X$ is a symmetric / Hermitian matrix.

\section*{1. Wigner Matrices}
We begin by fixing an infinite family of real-valued random variables $\left\{Y_{i j}\right\}_{j \geq i \geq 1}$. Then we can define a sequence of symmetric random matrices $\mathbf{Y}_{n}$ by

$$
\left[\mathbf{Y}_{n}\right]_{i j}=\left\{\begin{array}{ll}
Y_{i j}, & i \leq j \\
Y_{j i}, & i>j
\end{array} .\right.
$$

The matrix $\mathbf{Y}_{n}$ is symmetric and so has $n$ real eigenvalues, which we write in increasing order $\lambda_{1}\left(\mathbf{Y}_{n}\right) \leq \cdots \leq \lambda_{n}\left(\mathbf{Y}_{n}\right)$. In this very general setup, little can be said about these eigenvalues. The class of matrices we are going to begin studying, Wigner matrices, are given by the following conditions.

\begin{itemize}
  \item[-] We assume that the $\left\{Y_{i j}\right\}_{1 \leq i \leq j}$ are independent.
  \item[-] We assume that the diagonal entries $\left\{Y_{i i}\right\}_{i \geq 1}$ are identically-distributed, and the off-diagonal entries $\left\{Y_{i j}\right\}_{1 \leq i<j}$ are identically-distributed.
  \item[-] We assume that $\mathbb{E}\left(Y_{i j}^{2}\right)<\infty$ for all $i, j$. (I.e. $r_{2}=\max \left\{\mathbb{E}\left(Y_{11}^{2}\right), \mathbb{E}\left(Y_{12}^{2}\right)\right\}<\infty$.)
\end{itemize}

It is not just for convenience that we separate out the diagonal terms; as we will see, they really do not contribute to the eigenvalues in the limit as $n \rightarrow \infty$. It will also be convenient, at least at the start, to strengthen the final assumption to moments of all orders: for $k \geq 1$, set

$$
r_{k}=\max \left\{\mathbb{E}\left(Y_{11}^{k}\right), \mathbb{E}\left(Y_{12}^{k}\right)\right\} .
$$

To begin, we will assume that $r_{k}<\infty$ for each $k$; we will weaken this assumption later. Note: in the presence of higher moments, much of the following will not actually require identically-distributed entries. Rather, uniform bounds on all moments suffice. Herein, we will satisfy ourselves with the i.i.d. case.

One variation we will allow from i.i.d. is a uniform scaling of the matrices as $n$ increases. That is: let $\alpha_{n}$ be a sequence of positive scalars, and set

$$
\mathbf{X}_{n}=\alpha_{n} \mathbf{Y}_{n} .
$$

In fact, there is a natural choice for the scaling behavior, in order for there to be limiting behavior of eigenvalues. That is to say: we would like to arrange (if possible) that both $\lambda_{1}\left(\mathbf{X}_{n}\right)$ and $\lambda_{n}\left(\mathbf{X}_{n}\right)$ converge to finite numbers (different from each other). An easy way to test this is suggested by the following simple calculation:

$$
\lambda_{1}^{2}+\cdots+\lambda_{n}^{2} \leq n \cdot \max _{j=1}^{n}\left\{\lambda_{j}^{2}\right\}=n \cdot \max \left\{\lambda_{1}^{2}, \lambda_{n}^{2}\right\} .
$$

Hence, in order for $\lambda_{1}$ and $\lambda_{n}$ to both converge to distinct constants, it is necessary for the sequence $\frac{1}{n}\left(\lambda_{1}^{2}+\cdots+\lambda_{n}^{2}\right)$ to be bounded. Fortunately, this sequence can be calculated without explicit reference to eigenvalues: it is the (normalized) Hilbert-Schmidt norm (or Fröbenius norm) of a symmetric matrix X:

$$
\|\mathbf{X}\|_{2}^{2}=\frac{1}{n} \sum_{i=1}^{n} \lambda_{i}\left(\mathbf{X}_{n}\right)^{2}=\frac{1}{n} \sum_{1 \leq i, j \leq n} X_{i j}^{2} .
$$

Exercise 1.0.1. Verify the second equality above, by showing (using the spectral theorem) that both expressions are equal to the quantity $\frac{1}{n} \operatorname{Tr}\left(\mathbf{X}^{2}\right)$.

For our random matrix $\mathbf{X}_{n}$ above, then, we can calculate the expected value of this norm:

$$
\begin{aligned}
\mathbb{E}\left\|\mathbf{X}_{n}\right\|_{2}^{2}=\frac{1}{n} \alpha_{n}^{2} \sum_{1 \leq i, j \leq n} \mathbb{E}\left(Y_{i j}^{2}\right) & =\frac{\alpha_{n}^{2}}{n} \sum_{i=1}^{n} \mathbb{E}\left(Y_{i i}\right)^{2}+\frac{2 \alpha_{n}^{2}}{n} \sum_{1 \leq i<j \leq n} \mathbb{E}\left(Y_{i j}^{2}\right) \\
& =\alpha_{n}^{2} \cdot \mathbb{E}\left(Y_{11}^{2}\right)+(n-1) \alpha_{n}^{2} \cdot \mathbb{E}\left(Y_{12}^{2}\right) .
\end{aligned}
$$

We now have two cases. If $\mathbb{E}\left(Y_{12}^{2}\right)=0$ (meaning the off-diagonal terms are all 0 a.s.) then we see the "correct" scaling for $\alpha_{n}$ is $\alpha_{n} \sim 1$. This is a boring case: the matrices $\mathbf{X}_{n}$ are diagonal, with all diagonal entries identically distributed. Thus, these entries are also the eigenvalues, and so the distribution of eigenvalues is given by the common distribution of the diagonal entries. We ignore this case, and therefore assume that $\mathbb{E}\left(Y_{12}^{2}\right)>0$. Hence, in order for $\mathbb{E}\left\|\mathbf{X}_{n}\right\|_{2}^{2}$ to be a bounded sequence (that does not converge to 0 ), we must have $\alpha_{n} \sim n^{-1 / 2}$.\\
Definition 1.1. Let $\left\{Y_{i j}\right\}_{1 \leq i \leq j}$ and $\mathbf{Y}_{n}$ be as above, with $r_{2}<\infty$ and $\mathbb{E}\left(Y_{12}^{2}\right)>0$. Then the matrices $\mathbf{X}_{n}=n^{-1 / 2} \mathbf{Y}_{n}$ are Wigner matrices.

It is standard to abuse notation and refer to the sequence $\mathbf{X}_{n}$ as a a Wigner matrix. The preceding calculation shows that, if $\mathbf{X}_{n}$ is a Wigner matrix, then the expected Hilbert-Schmidt norm $\mathbb{E}\left\|\mathbf{X}_{n}\right\|_{2}^{2}$ converges (as $n \rightarrow \infty$ ) to the second moment of the (off-diagonal) entries. As explained above, this is prerequisite to the bulk convergence of the eigenvalues. As we will shortly see, it is also sufficient.

Consider the following example. Take the entires $Y_{i j}$ to be $N(0,1)$ normal random variables. These are easy to simulate with MATLAB. Figure 1 shows the histogram of all $n=4000$ eigenvalues of one instance of the corresponding Gaussian Wigner matrix $\mathbf{X}_{4000}$. The plot suggests that\\
$\lambda_{1}\left(\mathbf{X}_{n}\right) \rightarrow-2$ while $\lambda_{n}\left(\mathbf{X}_{n}\right) \rightarrow 2$ in this case. Moreover, although random fluctuations remain, it appears that the histogram of eigenvalues (sometimes called the density of eigenvalues) converges to a deterministic shape. In fact, this is universally true. There is a universal probability distribution $\sigma_{t}$ such that the density of eigenvalues of any Wigner matrix (with second moment $t$ ) converges to $\sigma_{t}$. The limiting distribution is known as Wigner's semicircle law:

$$
\sigma_{t}(d x)=\frac{1}{2 \pi t} \sqrt{\left(4 t-x^{2}\right)_{+}} d x .
$$

\begin{figure}[h]
\begin{center}
  \includegraphics[alt={},max width=\textwidth]{b29bb113-f2a1-4928-8b21-4b8b3f43679f-03_1139_1429_691_420}
\captionsetup{labelformat=empty}
\caption{Figure 1. The density of eigenvalues of an instance of $\mathbf{X}_{4000}$, a Gaussian Wigner matrix.}
\end{center}
\end{figure}

\section*{2. Wigner's Semicircle Law}
Theorem 2.1 (Wigner's Semicircle Law). Let $\mathbf{X}_{n}=n^{-1 / 2} \mathbf{Y}_{n}$ be a sequence of Wigner matrices, with entries satisfying $\mathbb{E}\left(Y_{i j}\right)=0$ for all $i, j$ and $\mathbb{E}\left(Y_{12}^{2}\right)=t$. Let $I \subset \mathbb{R}$ be an interval. Define the random variables

$$
E_{n}(I)=\frac{\#\left(\left\{\lambda_{1}\left(\mathbf{X}_{n}\right), \ldots, \lambda_{n}\left(\mathbf{X}_{n}\right)\right\} \cap I\right)}{n}
$$

Then $E_{n}(I) \rightarrow \sigma_{t}(I)$ in probability as $n \rightarrow \infty$.\\
The first part of this course is devoted to proving Wigner's Semicircle Law. The key observation (that Wigner made) is that one can study the behavior of the random variables $E_{n}(I)$ without computing the eigenvalues directly. This is accomplished by reinterpreting the theorem in terms of a random measure, the empirical law of eigenvalues.

Definition 2.2. Let $\mathbf{X}_{n}$ be a Wigner matrix. Its empirical law of eigenvalues $\mu_{\mathbf{X}_{n}}$ is the random discrete probability measure

$$
\mu_{\mathbf{X}_{n}}=\frac{1}{n} \sum_{j=1}^{n} \delta_{\lambda_{j}\left(\mathbf{X}_{n}\right)} .
$$

That is: $\mu_{\mathbf{X}_{n}}$ is defined as the (random) probability measure such that, for any continuous function $f \in C(\mathbb{R})$, the integral $\int f d \mu_{\mathbf{x}_{n}}$ is the random variable

$$
\int f d \mu_{\mathbf{X}_{n}}=\frac{1}{n} \sum_{j=1}^{n} f\left(\lambda_{j}\left(\mathbf{X}_{n}\right)\right) .
$$

Note that the random variables $E_{n}(I)$ in Theorem 2.1 are given by $E_{n}(I)=\int \mathbb{1}_{I} d \mu_{\mathbf{X}_{n}}$. Although $\mathbb{1}_{I}$ is not a continuous function, a simple approximation argument shows that the following theorem (which we also call Wigner's semicircle law) is a stronger version of Theorem 2.1.

Theorem 2.3 (Wigner's Semicircle Law). Let $\mathbf{X}_{n}=n^{-1 / 2} \mathbf{Y}_{n}$ be a sequence of Wigner matrices, with entries satisfying $\mathbb{E}\left(Y_{i j}\right)=0$ for all $i, j$ and $\mathbb{E}\left(Y_{12}^{2}\right)=t$. Then the empirical law of eigenvalues $\mu_{\mathbf{X}_{n}}$ converges in probability to $\sigma_{t}$ as $n \rightarrow \infty$. Precisely: for any $f \in C_{b}(\mathbb{R})$ (continuous bounded functions) and each $\epsilon>0$,

$$
\lim _{n \rightarrow \infty} \mathbb{P}\left(\left|\int f d \mu_{\mathbf{X}_{n}}-\int f d \sigma_{t}\right|>\epsilon\right)=0
$$

In this formulation, we can use the spectral theorem to eliminate the explicit appearance of eigenvalues in the law $\mu_{\mathbf{X}_{n}}$. Diagonalize $\mathbf{X}_{n}=U_{n}^{\top} \mathbf{\Lambda}_{n} U_{n}$. Then (by definition)

$$
\begin{aligned}
\int f d \mu_{\mathbf{X}_{n}}=\frac{1}{n} \sum_{j=1}^{n} f\left(\lambda_{j}\left(\mathbf{X}_{n}\right)\right) & =\frac{1}{n} \sum_{j=1}^{n} f\left(\left[\boldsymbol{\Lambda}_{n}\right]_{j j}\right) \\
& =\frac{1}{n} \operatorname{Tr} f\left(\boldsymbol{\Lambda}_{\mathbf{n}}\right)=\frac{1}{n} \operatorname{Tr}\left(U_{n}^{\top} f\left(\boldsymbol{\Lambda}_{\mathbf{n}}\right) U_{n}\right)=\frac{1}{n} \operatorname{Tr} f\left(\mathbf{X}_{n}\right) .
\end{aligned}
$$

The last equality is the statement of the spectral theorem. Usually we would use it in reverse to define $f\left(\mathbf{X}_{n}\right)$ for measurable $f$. However, in this case, we will take $f$ to be a polynomial. (Since both $\mu_{\mathbf{X}_{n}}$ and $\sigma_{t}$ are compactly-supported, any polynomial is equal to a $C_{b}(\mathbb{R})$ function on their supports, so this is consistent with Theorem 2.3.) This leads us to the third, a priori weaker form of Wigner's law.

Theorem 2.4 (Wigner's law for matrix moments). Let $\mathbf{X}_{n}=n^{-1 / 2} \mathbf{Y}_{n}$ be a sequence of Wigner matrices, with entries satisfying $\mathbb{E}\left(Y_{i j}\right)=0$ for all $i, j$ and $\mathbb{E}\left(Y_{12}^{2}\right)=t$. Let $f$ be a polynomial. Then the random variables $\int f d \mu_{\mathbf{x}_{n}}$ converge to $\int f d \sigma_{t}$ in probability as $n \rightarrow \infty$. Equivalently: for fixed $k \in \mathbb{N}$ and $\epsilon>0$,

$$
\lim _{n \rightarrow \infty} \mathbb{P}\left(\left|\frac{1}{n} \operatorname{Tr}\left(\mathbf{X}_{n}^{k}\right)-\int x^{k} \sigma_{t}(d x)\right|>\epsilon\right)=0
$$

Going back from Theorem 2.4 to Theorem 2.3 is, in principal, just a matter of approximating any $C_{b}$ function by polynomials on the supports of $\mu_{\mathbf{X}_{n}}$ and $\sigma_{t}$. However, even though $\mu_{\mathbf{X}_{n}}$ has finite support, there is no obvious bound on $\bigcup_{n} \operatorname{supp} \mu_{\mathbf{x}_{n}}$, and this makes such an approximation scheme tricky. In fact, when we eventually fully prove Theorem 2.3 (and hence Theorem 2.1 - approximating $\mathbb{1}_{I}$ by $C_{b}$ functions is routine in this case), we will take a different approach entirely. For now, however, it will be useful to have some intuition for the theorem, so we begin by presenting a scheme for proving Theorem 2.4. The place to begin is calculating the moments $\int x^{k} \sigma_{t}(d x)$.

Exercise 2.4.1. Let $\sigma_{t}$ be the semicircle law of variance $t$ defined above.

\begin{itemize}
  \item[(a)] With a suitable change of variables, show that $\int x^{k} \sigma_{t}(d x)=t^{k / 2} \int x^{k} \sigma_{1}(d x)$.
  \item[(b)] Let $m_{k}=\int x^{k} \sigma_{1}(d x)$. By symmetry, $m_{2 k+1}=0$ for all $k$. Use a trigonometric substitution to show that
$$
m_{0}=1, \quad m_{2 k}=\frac{2(2 k-1)}{k+1} m_{2(k-1)} .
$$
This recursion completely determines the even moments; show that, in fact,
$$
m_{2 k}=C_{k} \equiv \frac{1}{k+1}\binom{2 k}{k} .
$$
\end{itemize}

The numbers $C_{k}$ in Exercise 2.4.1 are the Catalan numbers. No one would disagree that they form the most important integer sequence in combinatorics. In fact, we will use some of their combinatorial interpretations to begin our study of Theorem 2.4.

\section*{3. Markov's Inequality and Convergence of Expectation}
To prove Theorem 2.4, we will begin by showing convergence of expectations. Once this is done, the result can be proved by showing the variance tends to 0, because of:

Lemma 3.1 (Markov's Inequality). Let $Y \geq 0$ be a random variable with $\mathbb{E}(Y)<\infty$. Then for any $\epsilon>0$,

$$
\mathbb{P}(Y>\epsilon) \leq \frac{\mathbb{E}(Y)}{\epsilon} .
$$

Proof. We simply note that $\mathbb{P}(Y>\epsilon)=\mathbb{E}\left(\mathbb{1}_{\{Y>\epsilon\}}\right)$. Now, on the event $\{Y>\epsilon\}, \frac{1}{\epsilon} Y>1$; on the complement of this event, $\mathbb{1}_{\{Y>\epsilon\}}=0$ while $\frac{1}{\epsilon} Y \geq 0$. Altogether, this means that we have the pointwise bound

$$
\frac{1}{\epsilon} Y \geq \mathbb{1}_{\{Y>\epsilon\}} .
$$

Taking expectations of both sides yields the result. \(\square\)

Corollary 3.2. Suppose $X_{n}$ is a sequence of random variables with $\mathbb{E}\left(X_{n}^{2}\right)<\infty$ for each $n$. Suppose that $\lim _{n \rightarrow \infty} \mathbb{E}\left(X_{n}\right)=m$, and that $\lim _{n \rightarrow \infty} \operatorname{Var}\left(X_{n}\right)=0$. Then $X_{n} \rightarrow m$ in probability: for each $\epsilon>0, \lim _{n \rightarrow \infty} \mathbb{P}\left(\left|X_{n}-m\right|>\epsilon\right)=0$.

Proof. Applying Markov's inequality to $Y=\left|X_{n}-m\right|$ (with $\epsilon^{2}$ in place of $\epsilon$ ),

$$
\mathbb{P}\left(\left|X_{n}-m\right|>\epsilon\right)=\mathbb{P}\left(\left(X_{n}-m\right)^{2}>\epsilon^{2}\right) \leq \frac{\mathbb{E}\left(\left(X_{n}-m\right)^{2}\right)}{\epsilon^{2}} .
$$

Now, using the triangle inequality for the $L^{2}(\mathbb{P})$-norm, we have

$$
\begin{aligned}
\mathbb{E}\left(\left(X_{n}-m\right)\right)^{2}=\left\|X_{n}-m\right\|_{L^{2}(\mathbb{P})}^{2} & =\left\|X_{n}-\mathbb{E}\left(X_{n}\right)+\mathbb{E}\left(X_{n}\right)-m\right\|_{L^{2}(\mathbb{P})}^{2} \\
& \left.\leq\left(\| X_{n}-\mathbb{E}\left(X_{n}\right)\right)\left\|_{L^{2}(\mathbb{P})}+\right\| \mathbb{E}\left(X_{n}\right)-m \|_{L^{2}(\mathbb{P})}\right)^{2} .
\end{aligned}
$$

Both $\mathbb{E}\left(X_{n}\right)$ and $m$ are constants, so the second $L^{2}(\mathbb{P})$ norm is simply $\left|\mathbb{E}\left(X_{n}\right)-m\right|$; the first term, on the other hand, is

$$
\left\|X_{n}-\mathbb{E}\left(X_{n}\right)\right\|_{L^{2}(\mathbb{P})}=\sqrt{\mathbb{E}\left[\left(X_{n}-\mathbb{E}\left(X_{n}\right)\right)^{2}\right]}=\sqrt{\operatorname{Var}\left(X_{n}\right)} .
$$

Hence, by Markov's inequality, we have

$$
\mathbb{P}\left(\left|X_{n}-m\right|>\epsilon\right) \leq \frac{\left(\sqrt{\operatorname{Var}\left(X_{n}\right)}+\left|\mathbb{E}\left(X_{n}\right)-m\right|\right)^{2}}{\epsilon^{2}} .
$$

The right-hand-side tends to 0 by assumption; this proves the result. \(\square\)

Remark 3.3. Using the Borel-Cantelli lemma, one can strengthen the result of Corollary 3.2: indeed, $X_{n} \rightarrow m$ almost surely, and so almost sure convergence also applies to Wigner's theorem. (This requires some additional estimates showing that the rate of convergence in probability is summably fast.) It is customary to state Wigner's theorem in terms of convergence in probability, though, because our setup (choosing all entries of all the matrices from one common probability space) is a little artificial. If the entries of $Y_{n}$ and $Y_{n+1}$ need not come from the same probability space, then almost sure convergence has no meaning; but convergence in probability still makes perfect sense.

Applying Corollary 3.2 to the terms in Theorem 2.4, this gives a two-step outline for a method of proof.

Step 1. For each odd $k$, show that $\frac{1}{n} \mathbb{E} \operatorname{Tr}\left(\mathbf{X}_{n}^{k}\right) \rightarrow 0$ as $n \rightarrow \infty$; for each even $k$, show that $\frac{1}{n} \mathbb{E} \operatorname{Tr}\left(\mathbf{X}_{n}^{k}\right) \rightarrow C_{k / 2}$ as $n \rightarrow \infty$.\\
Step 2. For each $k$, show that $\operatorname{Var}\left(\frac{1}{n} \operatorname{Tr}\left(\mathbf{X}_{n}^{k}\right)\right) \rightarrow 0$ as $n \rightarrow \infty$.\\
We will presently follow through with those two steps in the next section. The proof is quite combinatorial. We will also present a completely different, more analytical, proof of Theorem 2.3 in later lectures.

Remark 3.4. Note that the statement of Theorem 2.4 does not require higher moments of $\operatorname{Tr}\left(\mathbf{X}_{n}^{k}\right)$ to exist - the result is convergence in probability (and this is one benefit of stating the convergence this way). However: Step 1 above involves expectations, and so we will be taking expectations of terms that involve $k$ th powers of the variables $Y_{i j}$. Thus, in order to follow through with this program, we will need to make the additional assumption that $r_{k}<\infty$ for all $k$; the entries of $\mathbf{X}_{n}$ possess moments of all orders. This is a technical assumption that can be removed after the fact with appropriate cut-offs; when we delve into the analytical proof later, we will explain this.

\section*{4. First Proof of Wigner's Semicircle Law}
In this section, we give a complete proof of Theorem 2.4, under the assumption that all moments are finite.\\
4.1. Convergence of matrix moments in Expectation. We begin with Step 1: convergence in expectation.

Proposition 4.1. Let $\left\{Y_{i j}\right\}_{1 \leq i \leq j}$ be independent random variables, with $\left\{Y_{i i}\right\}_{i \geq 1}$ identically distributed and $\left\{Y_{i j}\right\}_{1 \leq i<j}$ identically distributed. Suppose that $r_{k}=\max \left\{\mathbb{E}\left(\left|Y_{11}\right|^{k}\right), \mathbb{E}\left(\left|Y_{12}\right|^{k}\right)\right\}<$ $\infty$ for each $k \in \mathbb{N}$. Suppose further than $\mathbb{E}\left(Y_{i j}\right)=0$ for all $i, j$ and set $t=\mathbb{E}\left(Y_{12}^{2}\right)$. If $i>j$, define $Y_{i j} \equiv Y_{j i}$, and let $\mathbf{Y}_{n}$ be the $n \times n$ matrix with $\left[\mathbf{Y}_{n}\right]_{i j}=Y_{i j}$ for $1 \leq i, j \leq n$. Let $\mathbf{X}_{n}=n^{-1 / 2} \mathbf{Y}_{n}$ be the corresponding Wigner matrix. Then

$$
\lim _{n \rightarrow \infty} \frac{1}{n} \mathbb{E} \operatorname{Tr}\left(\mathbf{X}_{n}^{k}\right)=\left\{\begin{array}{ll}
t^{k / 2} C_{k / 2}, & k \text { even } \\
0, & k \text { odd }
\end{array} .\right.
$$

Proof. To begin the proof, we expand the expected trace terms in terms of the entries. First, we have


\begin{equation*}
\frac{1}{n} \mathbb{E} \operatorname{Tr}\left(\mathbf{X}_{n}^{k}\right)=\frac{1}{n} \mathbb{E} \operatorname{Tr}\left[\left(n^{-1 / 2} \mathbf{Y}_{n}\right)^{k}\right]=n^{-k / 2-1} \mathbb{E} \operatorname{Tr}\left(\mathbf{Y}_{n}^{k}\right) . \tag{4.1}
\end{equation*}


Now, from the (repeated) definition of matrix multiplication, we have for any $1 \leq i, j \leq n$

$$
\left[\mathbf{Y}_{n}^{k}\right]_{i j}=\sum_{1 \leq i_{2}, \ldots, i_{k} \leq n} Y_{i i_{2}} Y_{i_{2} i_{3}} \cdots Y_{i_{k-1} i_{k}} Y_{i_{k} j} .
$$

Summing over the diagonal entries and taking the expectation, we get the expected trace:


\begin{equation*}
\mathbb{E} \operatorname{Tr}\left(\mathbf{Y}_{n}^{k}\right)=\sum_{i_{1}=1}^{n} \mathbb{E}\left(\left[\mathbf{Y}_{n}^{k}\right]_{i_{1} i_{1}}\right)=\sum_{1 \leq i_{1}, i_{2}, \ldots, i_{k} \leq n} \mathbb{E}\left(Y_{i_{1} i_{2}} Y_{i_{2} i_{3}} \cdots Y_{i_{k} i_{1}}\right) \equiv \sum_{\mathbf{i} \in[n]^{k}} \mathbb{E}\left(Y_{\mathbf{i}}\right), \tag{4.2}
\end{equation*}


where $[n]=\{1, \ldots, n\}$, and for $\mathbf{i}=\left(i_{1}, \ldots, i_{k}\right)$ we define $Y_{\mathbf{i}}=Y_{i_{1} i_{2}} \cdots Y_{i_{k} i_{1}}$. The term $\mathbb{E}\left(Y_{\mathbf{i}}\right)$ is determined by the $k$-index i, but only very weakly. The sum is better indexed by a collection of walks on graphs arising from the indices.

Definition 4.2. Let $\mathbf{i} \in[n]^{k}$ be a $k$-index, $\mathbf{i}=\left(i_{1}, i_{2}, \ldots, i_{k}\right)$. Define a graph $G_{\mathbf{i}}$ as follows: the vertices $V_{\mathbf{i}}$ are the distinct elements of $\left\{i_{1}, i_{2}, \ldots, i_{k}\right\}$, and the edges $E_{\mathbf{i}}$ are the distinct pairs among $\left\{i_{1}, i_{2}\right\},\left\{i_{2}, i_{3}\right\}, \ldots,\left\{i_{k-1}, i_{k}\right\},\left\{i_{k}, i_{1}\right\}$. The path $w_{\mathbf{i}}$ is the sequence

$$
w_{\mathbf{i}}=\left(\left\{i_{1}, i_{2}\right\},\left\{i_{2}, i_{3}\right\}, \ldots,\left\{i_{k-1}, i_{k}\right\},\left\{i_{k}, i_{1}\right\}\right)
$$

of edges.\\
For example, take $\mathbf{i}=(1,2,2,3,5,2,4,1,4,2)$; then

$$
V_{\mathbf{i}}=\{1,2,3,4,5\}, \quad E_{\mathbf{i}}=\{\{1,2\},\{1,4\},\{2,2\},\{2,3\},\{2,4\},\{2,5\},\{3,5\}\} .
$$

The index i which defines the graph now also defines a closed walk $w_{\mathbf{i}}$ on this graph. For this example, we have $Y_{\mathbf{i}}=Y_{12} Y_{22} Y_{23} Y_{35} Y_{52} Y_{24} Y_{41} Y_{14} Y_{42} Y_{21}$, which we can interpret as the walk $w_{\mathbf{i}}$ pictured below, in Figure 2. By definition, the walk $w_{\mathrm{i}}$ visits edge of $G_{\mathrm{i}}$ (including any self-edges present), beginning and ending at the beginning vertex. In particular, this means that the graph $G_{\mathrm{i}}$ is connected. The walk $w_{\mathrm{i}}$ encodes a labeling of the edges: the number of times each edge is traversed. We will denote this statistic also as $w_{\mathbf{i}}(e)$ for each edge $e$. In the above example, $w_{\mathbf{i}}(e)$ is equal to 2 on the edges $e \in\{\{1,2\},\{1,4\},\{2,4\}\}$, and is equal to 1 for all the other edges

\begin{figure}[h]
\begin{center}
  \includegraphics[alt={},max width=\textwidth]{b29bb113-f2a1-4928-8b21-4b8b3f43679f-09_707_534_217_802}
\captionsetup{labelformat=empty}
\caption{Figure 2. The graph and walk corresponding to the index $\mathbf{i}=(1,2,2,3,5,2,4,1,4,2)$.}
\end{center}
\end{figure}

(including the self-edge $\{2,2\}$ ). These numbers are actually evident in the expansion $Y_{\mathbf{i}}$ : using the fact that $Y_{i j}=Y_{j i}$, we have

$$
Y_{\mathbf{i}}=Y_{12}^{2} Y_{14}^{2} Y_{24}^{2} Y_{22} Y_{23} Y_{25} Y_{35},
$$

and in each case the exponent of $Y_{i j}$ is equal to $w_{\mathbf{i}}(\{i, j\})$. This is true in general (essentially by definition): if we define $w_{\mathbf{i}}(\{i, j\})=0$ if the pair $\{i, j\} \notin E_{\mathbf{i}}$, then


\begin{equation*}
Y_{\mathbf{i}}=\prod_{1 \leq i \leq j \leq n} Y_{i j}^{w_{\mathbf{i}}(\{i, j\})} . \tag{4.3}
\end{equation*}


Now, all the variables $Y_{i j}$ are independent. Since we allow the diagonal vs. off-diagonal terms to have different distributions, we should make a distinction between the self-edges $E_{\mathrm{i}}^{s}=\{\{i, i\} \in$ $\left.E_{\mathbf{i}}\right\}$ and connecting-edges $E_{\mathbf{i}}^{c}=\left\{\{i, j\} \in E_{\mathbf{i}}: i \neq j\right\}$. Then we have


\begin{equation*}
\mathbb{E}\left(Y_{\mathbf{i}}\right)=\prod_{1 \leq i \leq j \leq n} \mathbb{E}\left(Y_{i j}^{w_{\mathbf{i}}(\{i, j\})}\right)=\prod_{e_{s} \in E_{\mathbf{i}}^{s}} \mathbb{E}\left(Y_{11}^{w_{\mathbf{i}}\left(e_{s}\right)}\right) \cdot \prod_{e_{c} \in E_{\mathbf{i}}^{c}} \mathbb{E}\left(Y_{12}^{w_{\mathbf{i}}\left(e_{c}\right)}\right) . \tag{4.4}
\end{equation*}


That is: the value of $\mathbb{E}\left(Y_{\mathbf{i}}\right)$ is determined by the pair ( $G_{\mathbf{i}}, w_{\mathbf{i}}$ ) (or better yet $\left(E_{\mathbf{i}}, w_{\mathbf{i}}\right)$ ). Let us denote the common value in Equation 4.4 as $\Pi\left(G_{\mathbf{i}}, w_{\mathbf{i}}\right)$.

For any $k$-index $\mathbf{i}$, the connected oriented graph $G_{\mathbf{i}}$ has at most $k$ vertices. Since $w_{\mathbf{i}}$ records all the non-zero exponents in $Y_{\mathbf{i}}$ (that sum to the number of terms $k$ ), we have $\left|w_{\mathbf{i}}\right| \equiv \sum_{e \in E_{\mathbf{i}}} w_{\mathbf{i}}(e)=k$. Motivated by these conditions, let $\mathcal{G}_{k}$ denote the set of all pairs $(G, w)$ where $G=(V, E)$ is a connected graph with at most $k$ vertices, and $w$ is a closed walk covering $G$ satisfying $|w|=k$. Using Equation 4.4, we can reindex the sum of Equation 4.2 as


\begin{equation*}
\mathbb{E} \operatorname{Tr}\left(\mathbf{Y}_{n}^{k}\right)=\sum_{(G, w) \in \mathcal{G}_{k}} \sum_{\substack{\mathbf{i} \in[n]^{k} \\\left(G_{\mathbf{i}}, \mathbf{w}_{\mathbf{i}}\right)=(G, w)}} \mathbb{E}\left(Y_{\mathbf{i}}\right)=\sum_{(G, w) \in \mathcal{G}_{k}} \Pi(G, w) \cdot \#\left\{\mathbf{i} \in[n]^{k}:\left(G_{\mathbf{i}}, w_{\mathbf{i}}\right)=(G, w)\right\} . \tag{4.5}
\end{equation*}


Combining with the renormalization of Equation 4.1, we have


\begin{equation*}
\frac{1}{n} \mathbb{E} \operatorname{Tr}\left(\mathbf{X}_{n}^{k}\right)=\sum_{(G, w) \in \mathcal{G}_{k}} \Pi(G, w) \cdot \frac{\#\left\{\mathbf{i} \in[n]^{k}:\left(G_{\mathbf{i}}, w_{\mathbf{i}}\right)=(G, w)\right\}}{n^{k / 2+1}} . \tag{4.6}
\end{equation*}


It is, in fact, quite easy to count the sets of $k$-indices in Equation 4.6. For any $(G, w) \in \mathcal{G}_{k}$, an index with that corresponding graph $G$ and walk $w$ is completely determined by assigning which distinct values from $[n]$ appear at the vertices of $G$. For example, consider the pair $(G, w)$ in Figure 2. If $\mathbf{i}=\left(i_{1}, i_{2}, \ldots, i_{k}\right)$ is a $k$-index with this graph and walk, then reading along the walk we must have

$$
\mathbf{i}=\left(i_{1}, i_{2}, i_{2}, i_{3}, i_{5}, i_{2}, i_{4}, i_{1}, i_{4}, i_{2}, i_{1}\right)
$$

and so the set of all such $\mathbf{i}$ is determined by assigning 5 distinct values from $[n]$ to the indices $i_{1}, i_{2}, i_{3}, i_{4}, i_{5}$. There are $n(n-1)(n-2)(n-3)(n-4)$ ways of doing this. Following this example, in general we have:

Lemma 4.3. Given $(G, w) \in \mathcal{G}_{k}$, denote by $|G|$ the number of vertices in $G$. Then

$$
\#\left\{\mathbf{i} \in[n]^{k}:\left(G_{\mathbf{i}}, w_{\mathbf{i}}\right)=(G, w)\right\}=n(n-1) \cdots(n-|G|+1) .
$$

Hence, Equation 4.6 can be written simply as


\begin{equation*}
\frac{1}{n} \mathbb{E} \operatorname{Tr}\left(\mathbf{X}_{n}^{k}\right)=\sum_{(G, w) \in \mathcal{G}_{k}} \Pi(G, w) \cdot \frac{n(n-1) \cdots(n-|G|+1)}{n^{k / 2}+1} . \tag{4.7}
\end{equation*}


The summation is finite (since $k$ is fixed as $n \rightarrow \infty$ ), and so we are left to determined the values $\Pi(G, w)$. We begin with a simple observation: let $(G, w) \in \mathcal{G}_{k}$ and suppose there exists and edge $e=\{i, j\}$ with $w(e)=1$. This means that, in the expression 4.4 for $\Pi(G, w)$, a singleton term $\mathbb{E}\left(Y_{i j}^{w(e)}\right)=\mathbb{E}\left(Y_{i j}\right)$ appears. By assumption, the variables $Y_{i j}$ are all centered, and so this term is 0 ; hence, the product $\Pi(G, w)=0$ for any such pair $(G, w)$. This reduces the sum in Equation 4.6 considerably, since we need only consider those $w$ that cross each edge at least twice. We record this condition as $w \geq 2$, so


\begin{equation*}
\frac{1}{n} \mathbb{E} \operatorname{Tr}\left(\mathbf{X}_{n}^{k}\right)=\sum_{\substack{(G, w) \in \mathcal{G}_{k} \\ w \geq 2}} \Pi(G, w) \cdot \frac{n(n-1) \cdots(n-|G|+1)}{n^{k / 2+1}} . \tag{4.8}
\end{equation*}


The condition $w \geq 2$ restricts those graphs that can appear. Since $\left|w_{\mathbf{i}}\right|=k$, if each edge in $G_{\mathbf{i}}$ is traversed at least twice, this means that the number of edges is $\leq k / 2$.

Exercise 4.3.1. Let $G=(V, E)$ be a connected finite graph. Show that $|G|=\# V \leq \# E+1$, and that $|G|=\# V=\# E+1$ if and only if $G$ is a plane tree.

Hint: the claim is obvious for $\# V=1$. The inequality can be proved fairly easily by induction. The equality case also follows by studying this induction (if $G$ has a loop, removing the neighboring edges of a single vertex reduces the total number of edges by at least three).

Using Exercise 4.3.1, we see that for any graph $G=(V, E)$ appearing in the sum in Equation $4.8,|G| \leq k / 2+1$. The product $n(n-1) \cdots(n-|G|+1)$ is asymptotically equal to $n^{|G|}$, and so we see that the sequence $n \mapsto \frac{1}{n} \mathbb{E} \operatorname{Tr}\left(\mathbf{X}_{n}^{k}\right)$ is bounded. What's more, suppose that $k$ is odd. Since $|G| \leq \# E+1 \leq k / 2+1$ and $|G|$ is an integer, it follows that $|G| \leq(k-1) / 2+1=k / 2+1 / 2$. Hence, in this case, all the terms in the (finite $n$-independent) sum in Equation 4.8 are $O\left(n^{-1 / 2}\right)$. Thus, we have already proved that

$$
\lim _{n \rightarrow \infty} \mathbb{E} \operatorname{Tr}\left(\mathbf{X}_{n}^{k}\right)=0, \quad k \text { odd } .
$$

Henceforth, we assume that $k$ is even. In this case, it is still true that most of the terms in the sum in Equation 4.8 are 0. The following proposition testifies to this fact.

Proposition 4.4. Let $(G, w) \in \mathcal{G}_{k}$ with $w \geq 2$.

\begin{itemize}
  \item[(a)] If there exists a self-edge $e \in E_{s}$ in $G$, then $|G| \leq k / 2$.
  \item[(b)] If there exists an edge $e$ in $G$ with $w(e) \geq 3$, then $|G| \leq k / 2$.
\end{itemize}

Proof. (a) Since the graph $G=(V, E)$ contains a loop, it is not a tree; it follows from Exercise 4.3.1 that $\# V<\# E+1$. But $w \geq 2$ implies that $\# E \leq k / 2$, and so $\# V<k / 2+1$, and so $|G|=\# V \leq k / 2$.

\begin{itemize}
  \item[(b)] The sum of $w$ over all edges $E$ in $G$ is $k$. Hence, the sum of $w$ over $E \backslash\{e\}$ is $\leq k-3$. Since $w \geq 2$, this means that the number of edges excepting $e$ is $\leq(k-3) / 2$; hence, $\# E \leq(k-3) / 2+1=(k-1) / 2$. By the result of Exercise 4.3.1, this means that $\# V \leq(k-1) / 2+1=(k+1) / 2$. Since $k$ is even, it follows that $|G|=\# V \leq k / 2$. \(\square\)
\end{itemize}

Combining proposition 4.4 with Equation 4.8 suggest we drastically refine the set $\mathcal{G}_{k}$. Set $\mathcal{G}_{k}^{k / 2+1}$ to be the set of pairs $(G, w) \in \mathcal{G}_{k}$ where $G$ has $k / 2+1$ vertices, contains no self-edges, and the walk $w$ crosses every edge exactly 2 times. Then


\begin{equation*}
\frac{1}{n} \mathbb{E} \operatorname{Tr}\left(\mathbf{X}_{n}^{k}\right)=\sum_{(G, w) \in \mathcal{G}_{k}^{k / 2+1}} \Pi(G, w) \cdot \frac{n(n-1) \cdots(n-|G|+1)}{n^{k / 2+1}}+O_{k}\left(n^{-1}\right), \tag{4.9}
\end{equation*}


where $O_{k}\left(n^{-1}\right)$ means that the absolute difference of the terms is $\leq B_{k} / n$ for some $n$-independent constant $B_{k}$. Since $|G|=k / 2+1$ and $n(n-1) \cdots(n-k / 2+1) \sim n^{k / 2+1}$, it follows therefore that

$$
\lim _{n \rightarrow \infty} \mathbb{E} \operatorname{Tr}\left(\mathbf{X}_{n}^{k}\right)=\sum_{(G, w) \in \mathcal{G}_{k}^{k / 2+1}} \Pi(G, w) .
$$

Let $(G, w) \in \mathcal{G}_{k}^{k / 2+1}$. Since $w$ traverses each edge exactly twice, the number of edges in $G$ is $k / 2$. Since the number of vertices is $k / 2+1$, Exercise 4.3.1 shows that $G$ is a tree. In particular there are no self-edges (as we saw already in Proposition 4.4) and so the value of $\Pi(G, w)$ in Equation 4.4 is


\begin{equation*}
\Pi(G, w)=\prod_{e_{c} \in E^{c}} \mathbb{E}\left(Y_{12}^{w\left(e_{c}\right)}\right)=\prod_{e_{c} \in E^{c}} \mathbb{E}\left(Y_{12}^{2}\right)=t^{\# E}=t^{k / 2} . \tag{4.10}
\end{equation*}


Hence, we finally have


\begin{equation*}
\lim _{n \rightarrow \infty} \mathbb{E} \operatorname{Tr}\left(\mathbf{X}_{n}^{k}\right)=t^{k / 2} \cdot \# \mathcal{G}_{k}^{k / 2+1} . \tag{4.11}
\end{equation*}


Finally, we must enumerate the set $\mathcal{G}_{k}^{k / 2+1}$. To do so, we introduce another combinatorial structure: Dyck paths. Given a pair $(G, w) \in \mathcal{G}_{k}^{k / 2+1}$, define a sequence $\mathbf{d}=\mathbf{d}(G, w) \in\{+1,-1\}^{k}$ recursively as follows. Let $d_{1}=+1$. For $1<j \leq k$, if $w_{j} \notin\left\{w_{1}, \ldots, w_{j-1}\right\}$, set $d_{j}=+1$; otherwise, set $d_{j}=-1$; then $\mathbf{d}(G, w)=\left(d_{1}, \ldots, d_{k}\right)$. For example, with $k=10$ suppose that $G$ is the graph in Figure 3, with walk

$$
w=(\{1,2\},\{2,3\},\{3,2\},\{2,4\},\{4,5\},\{5,4\},\{4,6\},\{6,4\},\{4,2\},\{2,1\}) .
$$

Then $\mathbf{d}(G, w)=(+1,+1,-1,+1,+1,-1,+1,-1,-1,-1)$. One can interpret $\mathbf{d}(G, w)$ as a lattice path by looking at the successive sums: set $P_{0}=(0,0)$ and $P_{j}=\left(j, d_{1}+\cdots+d_{j}\right)$ for $1 \leq j \leq k$; then the piecewise linear path connecting $P_{0}, P_{1}, \ldots, P_{k}$ is a lattice path. Since $(G, w) \in \mathcal{G}_{k}^{2}$, each edge appears exactly two times in $w$, meaning that the ±1s come in pairs in $\mathbf{d}(G, w)$. Hence $d_{1}+\cdots+d_{k}=0$. What's more, for any edge $e$, the -1 assigned to its second

\begin{figure}[h]
\begin{center}
  \includegraphics[alt={},max width=\textwidth]{b29bb113-f2a1-4928-8b21-4b8b3f43679f-12_404_362_248_873}
\captionsetup{labelformat=empty}
\caption{Figure 3. A graph with associated walk in $\mathcal{G}_{10}^{2}$.}
\end{center}
\end{figure}

appearance in $w$ comes after the +1 corresponding to its first appearance; this means that the partial sums $d_{1}+\cdots+d_{j}$ are all $\geq 0$. That is: $\mathbf{d}(G, w)$ is a Dyck path.\\
Exercise 4.4.1. Let $k$ be even and let $\mathcal{D}_{k}$ denote the set of Dyck paths of length $k$

$$
\mathcal{D}_{k}=\left\{\left(d_{1}, \ldots, d_{k}\right) \in\{ \pm 1\}: \sum_{i=1}^{k} d_{i} \geq 0 \text { for } 1 \leq j \leq j \text {, and } \sum_{i=1}^{k} d_{i}=0\right\} .
$$

For $(G, w) \in \mathcal{G}_{k}^{k / 2+1}$, the above discussion shows that $\mathbf{d}(G, w) \in \mathcal{D}_{k}$. Show that $(G, w) \mapsto$ $\mathbf{d}(G, w)$ is a bijection $\mathcal{G}_{k}^{k / 2+1} \rightarrow \mathcal{D}_{k}$ by finding its explicit inverse.

It is well-known that $\# D_{k}=C_{k / 2}$ are enumerated by Catalan numbers. (For proof, see Stanley's Enumerative Combinatorics Vol. 2; or Lecture 7 from Math 247A Spring 2011, available at \href{http://www.math.ucsd.edu/}{www.math.ucsd.edu/} ${ }^{\sim}$ tkemp/247ASp11.) Combining this with Equation 4.10 concludes the proof. \(\square\)\\
4.2. Convergence of Matrix Moments in Probability. Now, to complete the proof of Theorem 2.4, we must now proceed with Step 2: we must show that $\frac{1}{n} \operatorname{Tr}\left(\mathbf{X}_{n}^{k}\right)$ has variance tending to 0 as $n \rightarrow \infty$. Using the ideas of the above proof, this is quite straightforward. We will actually show a little more: we will find the optimal rate that it tends to 0.

Proposition 4.5. Let $\left\{Y_{i j}\right\}_{1 \leq i \leq j}$ be independent random variables, with $\left\{Y_{i i}\right\}_{i \geq 1}$ identically distributed and $\left\{Y_{i j}\right\}_{1 \leq i<j}$ identically distributed. Suppose that $r_{k}=\max \left\{\mathbb{E}\left(\left|Y_{11}\right|^{k}\right), \mathbb{E}\left(\left|Y_{12}\right|^{k}\right)\right\}<$ $\infty$ for each $k \in \mathbb{N}$. Suppose further than $\mathbb{E}\left(Y_{i j}\right)=0$ for all $i, j$. If $i>j$, define $Y_{i j} \equiv Y_{j i}$, and let $\mathbf{Y}_{n}$ be the $n \times n$ matrix with $\left[\mathbf{Y}_{n}\right]_{i j}=Y_{i j}$ for $1 \leq i, j \leq n$. Let $\mathbf{X}_{n}=n^{-1 / 2} \mathbf{Y}_{n}$ be the corresponding Wigner matrix. Then

$$
\operatorname{Var}\left(\frac{1}{n} \operatorname{Tr}\left(\mathbf{X}_{n}^{k}\right)\right)=O_{k}\left(\frac{1}{n^{2}}\right) .
$$

Again, to be precise, saying $F(k, n)=O_{k}\left(1 / n^{2}\right)$ means that, for each $k$, there is a constant $B_{k}<\infty$ so that $F(k, n) \leq B_{k} / n^{2}$.

Proof. We proceed as above, expanding the variance in terms of the entries of the matrix.

$$
\begin{aligned}
\operatorname{Var}\left(\frac{1}{n} \operatorname{Tr}\left(\mathbf{X}_{n}^{k}\right)\right) & =\mathbb{E}\left[\frac{1}{n} \cdot \frac{1}{n^{k / 2}} \operatorname{Tr}\left(\mathbf{Y}_{n}^{k}\right)\right]^{2}-\left(\mathbb{E}\left[\frac{1}{n} \cdot \frac{1}{n^{k / 2}} \operatorname{Tr}\left(\mathbf{Y}_{n}^{k}\right)\right]\right)^{2} \\
& =\frac{1}{n^{k+2}}\left\{\mathbb{E}\left[\operatorname{Tr}\left(\mathbf{Y}_{n}^{k}\right)\right]^{2}-\left(\mathbb{E} \operatorname{Tr}\left(\mathbf{Y}_{n}^{k}\right)\right)^{2}\right\}
\end{aligned}
$$

Expanding the trace exactly as above, and squaring, gives

$$
\mathbb{E}\left[\operatorname{Tr}\left(\mathbf{Y}_{n}^{k}\right)\right]^{2}=\sum_{\mathbf{i}, \mathbf{j} \in[n]^{k}} \mathbb{E}\left(Y_{\mathbf{i}} Y_{\mathbf{j}}\right), \quad\left(\mathbb{E} \operatorname{Tr}\left(\mathbf{Y}_{n}^{k}\right)\right)^{2}=\sum_{\mathbf{i}, \mathbf{j} \in[n]^{k}} \mathbb{E}\left(Y_{\mathbf{i}}\right) \mathbb{E}\left(Y_{\mathbf{j}}\right) .
$$

Thus the variance in question is

$$
\operatorname{Var}\left(\frac{1}{n} \operatorname{Tr}\left(\mathbf{X}_{n}^{k}\right)\right)=\frac{1}{n^{k+2}} \sum_{\mathbf{i}, \mathbf{j} \in[n]^{k}}\left[\mathbb{E}\left(Y_{\mathbf{i}} Y_{\mathbf{j}}\right)-\mathbb{E}\left(Y_{\mathbf{i}}\right) \mathbb{E}\left(Y_{\mathbf{j}}\right)\right] .
$$

Now, as above, the values $\mathbb{E}\left(Y_{\mathbf{i}} Y_{\mathbf{j}}\right)$ and $\mathbb{E}\left(Y_{\mathbf{i}}\right) \mathbb{E}\left(Y_{\mathbf{j}}\right)$ only depend on $\mathbf{i}$ and $\mathbf{j}$ through a certain graph structure underlying the $2 k$-tuple $\mathbf{i}, \mathbf{j}$. Indeed, let $\mathbf{G}_{\mathbf{i} \# \mathbf{j}} \equiv \mathbf{G}_{\mathbf{i}} \cup \mathbf{G}_{\mathbf{j}}$, where the union of two graphs (whose vertices are chosen from the same underlying set, in this case $[n]$ ) is the graph whose vertex set is the union of the vertex sets, and whose edge set is the union of the edge sets. For example, if $\mathbf{i}=(1,2,1,3,2,3)$ and $\mathbf{j}=(2,4,5,1,4,3)$, then

$$
Y_{\mathbf{i}} Y_{\mathbf{j}}=Y_{12} Y_{21} Y_{13} Y_{32} Y_{23} Y_{31} \cdot Y_{24} Y_{45} Y_{51} Y_{14} Y_{43} Y_{32}
$$

and so

$$
G_{\mathbf{i} \# \mathbf{j}}=(\{1,2,3,4,5\},\{\{1,2\},\{1,3\},\{2,3\},\{2,4\},\{4,5\},\{1,5\},\{1,4\},\{3,4\}\}) .
$$

The word $Y_{\mathbf{i}} Y_{\mathbf{j}}$ also gives rise to the two graph walks $w_{\mathbf{i}}$ and $w_{\mathbf{j}}$; as the above example shows, they do not constitute a single walk of length $2 k$, since there is no a priori reason for the endpoint of the first walk to coincide with the starting point of the second walk.

Since we can recover $\left(G_{\mathbf{i}}, w_{\mathbf{i}}\right)$ and $\left(G_{\mathbf{j}}, w_{\mathbf{j}}\right)$ from $\left(G_{\mathbf{i} \# \mathbf{j}}, w_{\mathbf{i}}, w_{\mathbf{j}}\right)$, the same reasoning as in the previous proposition shows that the quantity $\mathbb{E}\left(Y_{\mathbf{i}} Y_{\mathbf{j}}\right)-\mathbb{E}\left(Y_{\mathbf{i}}\right) \mathbb{E}\left(Y_{\mathbf{j}}\right)$ is determined by the data $\left(G_{\mathbf{i} \# \mathbf{j}}, w_{\mathbf{i}}, w_{\mathbf{j}}\right)$. Denote this common value as

$$
\mathbb{E}\left(Y_{\mathbf{i}} Y_{\mathbf{j}}\right)-\mathbb{E}\left(Y_{\mathbf{i}}\right) \mathbb{E}\left(Y_{\mathbf{j}}\right)=\pi\left(G_{\mathbf{i} \# \mathbf{j}}, w_{\mathbf{i}}, w_{\mathbf{j}}\right) .
$$

Let $\mathcal{G}_{k, k}$ be the set of connected graphs $G$ with $\leq 2 k$ vertices, together with two paths each of length $k$ whose union covers $G$. This means we can expand (as above)

$$
\operatorname{Var}\left(\frac{1}{n} \operatorname{Tr}\left(\mathbf{X}_{n}^{k}\right)\right)=\frac{1}{n^{k+2}} \sum_{\left(G, w, w^{\prime}\right) \in \mathcal{G}_{k, k}} \pi\left(G, w, w^{\prime}\right) \cdot \#\left\{(\mathbf{i}, \mathbf{j}) \in[n]^{2 k}:\left(G_{\mathbf{i} \# \mathbf{j}}, w_{\mathbf{i}}, w_{\mathbf{j}}\right)=\left(G, w, w^{\prime}\right)\right\} .
$$

As before, let $E_{\mathbf{i} \# \mathbf{j}}^{s}$ denote the self edges and $E_{\mathbf{i} \# \mathbf{j}}^{c}$ the connecting edges in $G_{\mathbf{i} \# \mathbf{j}}$. For any edge $e$ in $G_{\mathbf{i} \# \mathbf{j}}$, let $w_{\mathbf{i} \# \mathbf{j}}(e)$ denote the number of times the edge $e$ is traversed by either of the two paths $w_{\mathbf{i}}$ and $w_{\mathbf{j}}$. Then, as in (4.4),

$$
\begin{aligned}
\mathbb{E}\left(Y_{\mathbf{i}} Y_{\mathbf{j}}\right)-\mathbb{E}\left(Y_{\mathbf{i}}\right) \mathbb{E}\left(Y_{\mathbf{j}}\right) & =\prod_{e_{s} \in E_{\mathbf{i} \# \mathbf{j}}} \mathbb{E}\left(Y_{11}^{w_{i \neq \mathbf{j}}\left(e_{s}\right)}\right) \cdot \prod_{e_{c} \in E_{\mathbf{i} \# \mathbf{j}}^{c}} \mathbb{E}\left(Y_{12}^{w_{\mathbf{i} \# \mathbf{j}}\left(e_{c}\right)}\right) \\
& -\prod_{e_{s} \in E_{\mathbf{i}}^{\mathbf{s}}} \mathbb{E}\left(Y_{11}^{w_{\mathbf{i}}\left(e_{s}\right)}\right) \cdot \prod_{e_{c} \in E_{\mathbf{i}}^{c}} \mathbb{E}\left(Y_{12}^{w_{\mathbf{i}}\left(e_{c}\right)}\right) \cdot \prod_{e_{s} \in E_{\mathbf{j}}^{\mathrm{s}}} \mathbb{E}\left(Y_{11}^{w_{\mathbf{j}}\left(e_{s}\right)}\right) \cdot \prod_{e_{c} \in E_{\mathbf{j}}^{c}} \mathbb{E}\left(Y_{12}^{w_{\mathbf{j}}\left(e_{c}\right)}\right) .
\end{aligned}
$$

The key thing is to note that

$$
\sum_{e \in G_{\mathbf{i} \# \mathbf{j}}} w_{\mathbf{i} \# \mathbf{j}}(e)=2 k=\sum_{e \in G_{\mathbf{i}}} w_{\mathbf{i}}(e)+\sum_{e \in G_{\mathbf{j}}} w_{\mathbf{j}}(e) ;
$$

that is, the sum of all exponents in either term $\mathbb{E}\left(Y_{\mathbf{i}} Y_{\mathbf{j}}\right)$ or $\mathbb{E}\left(Y_{\mathbf{i}}\right) \mathbb{E}\left(Y_{\mathbf{j}}\right)$ is the total length of the words, which is $2 k$. Thus, each of these terms is bounded (in modulus) above by something of the form $r_{m_{1}} \cdots r_{m_{\ell}}$ for some positive integers $m_{1}, \ldots m_{\ell}$ for which $m_{1}+\cdots+m_{\ell}=2 k$. (Recall that\\
$r_{m}=\max \left\{\mathbb{E}\left(\left|Y_{11}\right|^{m}\right), \mathbb{E}\left(\left|Y_{12}\right|^{m}\right)\right\}$.) There are only finitely many such integer partitions of $2 k$, and so the maximum $M_{2 k}$ over all of them is finite. We therefore have the (blunt) upper bound

$$
\left|\pi\left(G_{\mathbf{i} \# \mathbf{j}}, w_{\mathbf{i}}, w_{\mathbf{j}}\right)\right|=\left|\mathbb{E}\left(Y_{\mathbf{i}} Y_{\mathbf{j}}\right)-\mathbb{E}\left(Y_{\mathbf{i}}\right) \mathbb{E}\left(Y_{\mathbf{j}}\right)\right| \leq 2 M_{2 k}, \quad \forall \mathbf{i}, \mathbf{j} \in[n]^{k} .
$$

That being said, we can (as in the previous proposition) show that many of these terms are identically 0. Indeed, following the reasoning from above: by construction, every edge in the joined graph $G_{\mathrm{i} \# \mathrm{j}}$ is traversed at least once by the union of the two paths $w_{\mathrm{i}}$ and $w_{\mathrm{j}}$. Suppose that $e$ is an edge that is traversed only once. This means that $w_{\mathbf{i} \# \mathbf{j}}(e)=1$, and so it follows that the two values $w_{\mathbf{i}}(e), w_{\mathbf{j}}(e)$ are $\{0,1\}$. Hence, the above expansion and the fact that $\mathbb{E}\left(Y_{11}\right)=\mathbb{E}\left(Y_{12}\right)=0$ show that $\pi\left(G_{\mathbf{i} \# \mathbf{j}}, w_{\mathbf{i}}, w_{\mathbf{j}}\right)=0$ in this case. If we think if a path as a labeling of edges (counting the number of times an edge is traversed), then this means the variance sum reduces to

$$
\operatorname{Var}\left(\frac{1}{n} \operatorname{Tr}\left(\mathbf{X}_{n}^{k}\right)\right)=\sum_{\substack{\left(G, w, w^{\prime}\right) \in \mathcal{G}_{k, k} \\ w+w^{\prime} \geq 2}} \pi\left(G, w, w^{\prime}\right) \cdot \frac{\#\left\{(\mathbf{i}, \mathbf{j}) \in[n]^{2 k}:\left(G_{\mathbf{i} \# \mathbf{j}}, w_{\mathbf{i}}, w_{\mathbf{j}}\right)=\left(G, w, w^{\prime}\right)\right\}}{n^{k+2}} .
$$

The enumeration of the number of $2 k$-tuples yielding a certain graph with two walks is the same as in the previous proposition: the structure ( $G, w, w^{\prime}$ ) specifies the $2 k$-tuple precisely once we select the $|G|$ distinct indices for the vertices. So, as before, this ratio becomes

$$
\frac{n(n-1) \cdots(n-|G|+1)}{n^{k+2}} .
$$

Now, we have the condition $w+w^{\prime} \geq 2$, meaning every edge is traversed at least twice. Since there are $k$ steps in each of the two paths, this means there are at most $k$ edges. Appealing again to Exercise 4.3.1, it follows that $|G| \leq k+1$. Hence, $n(n-1) \cdots(n-|G|+1) \leq n^{|G|} \leq n^{k+1}$, and so we have already proved that

$$
\operatorname{Var}\left(\frac{1}{n} \operatorname{Tr}\left(\mathbf{X}_{n}^{k}\right)\right) \leq \sum_{\substack{\left(G, w, w^{\prime}\right) \in \mathcal{G}_{k, k} \\ w+w^{\prime} \geq 2}} \pi\left(G, w, w^{\prime}\right) \cdot \frac{n^{k+1}}{n^{k+2}} \leq \frac{1}{n} \cdot 2 M_{2 k} \cdot \# \mathcal{G}_{k, k} .
$$

The (potentially enormous) number $B_{k}=2 M_{2 k} \cdot \# \mathcal{G}_{k, k}$ is independent of $n$, and so we have already proved that the variance is $=O_{k}(1 / n)$, which is enough to prove Theorem 2.4 (convergence in probability).

We will go one step further and show that there are also no terms with $|G|=k+1$. From Exercise 4.3.1 again, since there are at most $k$ edges, it follows that there are exactly $k$ edges and so $G$ is a tree; this means that $w+w^{\prime}=2$, every edge is traversed exactly two times. It then follows that $G_{\mathrm{i}}$ and $G_{\mathrm{j}}$ share no edges in common. Indeed, suppose $e$ is a shared edge. Each walk $w_{\mathbf{i}}, w_{\mathbf{j}}$ traverses all edges of its graph $G_{\mathbf{i}}, G_{\mathbf{j}}$. Since $w_{\mathbf{i}}(e)+w_{\mathbf{j}}(e)=2$, the values $\left(w_{\mathbf{i}}(e), w_{\mathbf{j}}(e)\right)$ are either $(2,0),(1,1)$, or (0, 2). The first and last cases are impossible: in the first case, this would mean that $w_{\mathbf{j}}$ does not traverse $e$, which is an edge in $G_{\mathbf{j}}$, contradicting the definition of $w_{\mathbf{j}}$. But $(1,1)$ is also impossible: since the union graph $G_{\mathbf{i} \# \mathbf{j}}$ is a tree, each subgraph $G_{\mathbf{i}}$ and $G_{\mathbf{j}}$ is a tree, and since the walks $w_{\mathrm{i}}$ and $w_{\mathrm{j}}$ cover each edge and return to their starting points, each edge must be traversed an even number of times (as there are no loops).

Hence, we see that the only graph walks $\left(G, w, w^{\prime}\right) \in \mathcal{G}_{k, k}$ with $|G|=k+1$ must have the edge sets covered by $w$ and $w^{\prime}$ distinct. In other words, if $\left(G_{\mathbf{i} \# \mathbf{j}}, w_{\mathbf{i}}, w_{\mathbf{j}}\right)=\left(G, w, w^{\prime}\right)$, then the edge sets do not intersect: $\left\{\left\{i_{1}, i_{2}\right\}, \ldots,\left\{i_{k}, i_{1}\right\}\right\} \cap\left\{\left\{j_{1}, j_{2}\right\}, \ldots,\left\{j_{k}, j_{1}\right\}\right\}=\varnothing$. But that means that the products $Y_{\mathbf{i}}$ and $Y_{\mathbf{j}}$ are independent, and so $\pi\left(G_{\mathbf{i} \# \mathbf{j}}, w_{\mathbf{i}}, w_{\mathbf{j}}\right)=\mathbb{E}\left(Y_{\mathbf{i}} Y_{\mathbf{j}}\right)-\mathbb{E}\left(Y_{\mathbf{i}}\right) \mathbb{E}\left(Y_{\mathbf{j}}\right)=0$. Thus, we\\
have shown that, for any $\left(G, w, w^{\prime}\right) \in \mathcal{G}_{k, k}$ with $k+1$ vertices, $\pi\left(G, w, w^{\prime}\right)=0$. So, we actually have

$$
\operatorname{Var}\left(\frac{1}{n} \operatorname{Tr}\left(\mathbf{X}_{n}^{k}\right)\right)=\sum_{\substack{\left(G, w, w^{\prime}\right) \in \mathcal{G}_{k, k} \\ w+w^{\prime} \geq 2,[G] \leq k}} \pi\left(G, w, w^{\prime}\right) \cdot \frac{\#\left\{(\mathbf{i}, \mathbf{j}) \in[n]^{2 k}:\left(G_{\mathbf{i} \# \mathbf{j}}, w_{\mathbf{i}}, w_{\mathbf{j}}\right)=\left(G, w, w^{\prime}\right)\right\}}{n^{k+2}},
$$

and so the same counting argument now bounded the variance by

$$
\operatorname{Var}\left(\frac{1}{n} \operatorname{Tr}\left(\mathbf{X}_{n}^{k}\right)\right) \leq \sum_{\substack{\left(G, w, w^{\prime}\right) \in \mathcal{G}_{k, k} \\ w+w^{\prime} \geq 2,|G| \leq k}} \pi\left(G, w, w^{\prime}\right) \cdot \frac{n^{k}}{n^{k+2}} \leq \frac{1}{n^{2}} \cdot 2 M_{2 k} \cdot \# \mathcal{G}_{k, k} .
$$

This proves the proposition. \(\square\)

Remark 4.6. We showed not only that $\operatorname{Var}\left[\frac{1}{n} \operatorname{Tr}\left(\mathbf{X}_{n}^{k}\right)\right] \rightarrow 0$ as $n \rightarrow \infty$, giving convergence in probability by Chebyshev's inequality; we actually proved that the variance is $O\left(1 / n^{2}\right)$. This is relevant because $\sum_{n} \frac{1}{n^{2}}<\infty$. It therefore follows from the Borel-Cantelli lemma that

$$
\frac{1}{n} \operatorname{Tr}\left(\mathbf{X}_{n}^{k}\right) \rightarrow \int x^{k} \sigma_{t}(d x) \text { a.s. }
$$

as $n \rightarrow \infty$, with the caveat that this almost sure convergence only makes sense if we artificially sample all entries of all matrices from the same probability space.\\
4.3. Weak Convergence. We have now proved Theorem 2.4, which was a weaker form of Theorem 2.3. In fact, we can now fairly easily prove this stronger theorem. Let us first remark that we may easily fix the variance of the entries $t$ to be $t=1$ : an elementary scaling argument then extends Theorem 2.3 to the general case. With this convention in hand, we begin with the following useful cutoff lemma.

Lemma 4.7. Let $k \in \mathbb{N}$ and $\epsilon>0$. Then for any $b>4$,

$$
\limsup _{n \rightarrow \infty} \mathbb{P}\left(\int_{|x|>b}|x|^{k} \mu_{\mathbf{X}_{n}}(d x)>\epsilon\right)=0 .
$$

Proof. First, by Markov's inequality, we have

$$
\mathbb{P}\left(\int_{|x|>b}|x|^{k} \mu_{\mathbf{X}_{n}}(d x)>\epsilon\right) \leq \frac{1}{\epsilon} \mathbb{E}\left(\int_{|x|>b}|x|^{k} \mu_{\mathbf{X}_{n}}(d x)\right) .
$$

Now, let $\nu$ be the random measure $\nu(d x)=|x|^{k} \mu_{\mathbf{X}_{n}}(d x)$. Markov's inequality (which applies to all positive measures, not just probability measures) now shows that

$$
\int_{|x|>b}|x|^{k} \mu_{\mathbf{X}_{n}}(d x)=\nu\{x:|x|>b\}=\nu\left\{x:|x|^{k}>b^{k}\right\} \leq \frac{1}{b^{k}} \int|x|^{k} \nu(d x)=\frac{1}{b^{k}} \int x^{2 k} \mu_{\mathbf{X}_{n}}(d x) .
$$

So, taking expectations, we have

$$
\mathbb{P}\left(\int_{|x|>b}|x|^{k} \mu_{\mathbf{X}_{n}}(d x)>\epsilon\right) \leq \frac{1}{\epsilon b^{k}} \mathbb{E}\left(\int x^{2 k} \mu_{\mathbf{X}_{n}}(d x)\right)=\frac{1}{\epsilon b^{k}} \cdot \frac{1}{n} \mathbb{E} \operatorname{Tr}\left(\mathbf{X}_{n}^{k}\right) .
$$

Now, by Proposition 4.1, the right hand side converges to $C_{k} / \epsilon b^{k}$ where $C_{k}$ is the Catalan number, which is bounded by $4^{k}$. Hence, it follows that

$$
\limsup _{n \rightarrow \infty} \mathbb{P}\left(\int_{|x|>b}|x|^{k} \mu_{\mathbf{X}_{n}}(d x)>\epsilon\right) \leq \frac{1}{\epsilon}\left(\frac{4}{b}\right)^{k} .
$$

On the other hand, when $|x|>b>4>1$, the function $k \mapsto|x|^{k}$ is strictly increasing, which means that the sequence of lim sups on the left-hand-side is increasing. But this sequence decays exponentially since $4 / b<1$. The only way this is possible is if the sequence of lim sups is constantly 0, as claimed. \(\square\)

Proof of Theorem 2.3. Fix a bounded continuous function $f \in C_{b}(\mathbb{R})$, fix $\epsilon>0$, and fix $b>4$. By the Weierstrass approximation theorem, there is a polynomial $P_{\epsilon}$ such that

$$
\sup _{|x| \leq b}\left|f(x)-P_{\epsilon}(x)\right|<\frac{\epsilon}{6} .
$$

Now, we have the triangle inequality estimates

$$
\begin{aligned}
\left|\int f d \mu_{\mathbf{X}_{n}}-\int f d \sigma_{1}\right| & \leq\left|\int f d \mu_{\mathbf{X}_{n}}-\int P_{\epsilon} d \mu_{\mathbf{X}_{n}}\right|+\left|\int P_{\epsilon} d \mu_{\mathbf{X}_{n}}-\int P_{\epsilon} d \sigma_{1}\right| \\
& +\left|\int P_{\epsilon} d \sigma_{1}-\int f d \sigma_{1}\right| .
\end{aligned}
$$

Hence, the event $\left\{\left|\int f d \mu_{\mathbf{X}_{n}}-\int f d \sigma_{1}\right|>\epsilon\right\}$ is contained in the union of the three events that each of the four above terms is $>\epsilon / 3$. But this means that

$$
\begin{aligned}
\mathbb{P}\left(\left|\int f d \mu_{\mathbf{X}_{n}}-\int f d \sigma_{1}\right|>\epsilon\right) & \leq \mathbb{P}\left(\left|\int f d \mu_{\mathbf{X}_{n}}-\int P_{\epsilon} d \mu_{\mathbf{X}_{n}}\right|>\epsilon / 3\right) \\
& +\mathbb{P}\left(\left|\int P_{\epsilon} d \mu_{\mathbf{X}_{n}}-\int P_{\epsilon} d \sigma_{1}\right|>\epsilon / 3\right) \\
& +\mathbb{P}\left(\left|\int P_{\epsilon} d \sigma_{1}-\int f d \sigma_{1}\right|>\epsilon / 3\right) .
\end{aligned}
$$

By construction, $\left|P_{\epsilon}-f\right|<\epsilon / 6$ on $[-b, b]$, which includes the support $[-2,2]$ of $\sigma_{1}$; thus, the last term is identically 0. For the first term, we break up the integral over $[-b, b]$ and its complement:

$$
\left|\int\left(f-P_{\epsilon}\right) d \mu_{\mathbf{X}_{n}}\right| \leq \int_{|x| \leq b}\left|f(x)-P_{\epsilon}(x)\right| \mu_{\mathbf{X}_{n}}(d x)+\int_{|x|>b}\left|f(x)-P_{\epsilon}(x)\right| \mu_{\mathbf{X}_{n}}(d x) .
$$

By the same reasoning as above, we can estimate

$$
\begin{aligned}
\mathbb{P}\left(\left|\int f d \mu_{\mathbf{X}_{n}}-\int P_{\epsilon} d \mu_{\mathbf{X}_{n}}\right|>\epsilon / 3\right) & \leq \mathbb{P}\left(\int\left|f-P_{\epsilon}\right| \mathbb{1}_{|x| \leq b} d \mu_{\mathbf{X}_{n}}>\epsilon / 6\right) \\
& +\mathbb{P}\left(\int\left|f-P_{\epsilon}\right| \mathbb{1}_{|x|>b} d \mu_{\mathbf{X}_{n}}>\epsilon / 6\right) .
\end{aligned}
$$

Again, by construction, $\left|f-P_{\epsilon}\right|<\epsilon / 6$ on $[-b, b]$, and so since $\mu_{\mathbf{x}_{n}}$ is a probability measure, the first term is identically 0. We therefore have the estimate


\begin{align*}
\mathbb{P}\left(\left|\int f d \mu_{\mathbf{X}_{n}}-\int f d \sigma_{1}\right|>\epsilon\right) & \leq \mathbb{P}\left(\int\left|f-P_{\epsilon}\right| \mathbb{1}_{|x|>b} d \mu_{\mathbf{X}_{n}}>\epsilon / 6\right)  \tag{4.12}\\
& +\mathbb{P}\left(\left|\int P_{\epsilon} d \mu_{\mathbf{X}_{n}}-\int P_{\epsilon} d \sigma_{1}\right|>\epsilon / 3\right) . \tag{4.13}
\end{align*}


That the terms in (4.13) tend to 0 as $n \rightarrow \infty$ follows immediately from Theorem 2.4 (since convergence in probability respects addition and scalar multiplication). So we are left only to estimate (4.12). We do this as follows. Let $k=\operatorname{deg} P_{\epsilon}$. Since $f$ is bounded, we have $\mid f(x)-$\\
$P_{\epsilon}(x)\left|\leq\|f\|_{\infty}+\left|P_{\epsilon}(x)\right|\right.$, and on the set $| x \mid>b$, this is $\leq c|x|^{k}$ for some constant $c$ (since $b>0$ ). This means that

$$
\mathbb{P}\left(\int\left|f-P_{\epsilon}\right| 1_{|x| \geq b} d \mu_{\mathbf{X}_{n}}>\epsilon / 6\right) \leq \mathbb{P}\left(\int c|x|^{k} \mathbb{1}_{|x| \geq b} \mu_{\mathbf{X}_{n}}(d x)>\epsilon / 6\right),
$$

and therefore, by Lemma 4.7, this sequence has $\limsup _{n \rightarrow \infty}=0$. This concludes the proof. \(\square\)

Remark 4.8. From (4.13) and our proof of Theorem 2.4, we actually have decay rate $O_{k}\left(1 / n^{2}\right)$ for (4.13). A much more careful analysis of the proof of Lemma 4.7 shows that (4.12) also decays summably fast, and so we can actually conclude almost sure weak convergence in general. We will reprove this result in a different way later in these notes.

Finally, this brings us to the proof of Theorem 2.1.\\
Proof of Theorem 2.1. The random variable $E_{n}(I)$ in the statement of the theorem can be written as

$$
E_{n}(I)=\frac{1}{n} \cdot \#\left\{j \in[n]: \lambda_{j}\left(\mathbf{X}_{n}\right) \in I\right\}=\mu_{\mathbf{X}_{n}}(I) .
$$

That is to say, the desired conclusion is that $\mu_{\mathbf{X}_{n}}(I) \rightarrow \sigma_{t}(I)$ in probability for all intervals $I$. By an easy scaling argument, we can take $t=1$, in which case we have proved above that $\mu_{\mathbf{x}_{n}} \rightarrow \sigma_{1}$ weakly in probability. By a standard convergence theorem in probability theory (c.f. Theorem 25.8 in Billingsley's "Probability and Measure" (Third Edition)), this implies that $\mu_{\mathbf{X}_{n}}(A) \rightarrow \sigma_{1}(A)$ in probability for every $\sigma_{1}$-continuous measurable seat $A$. But since $\sigma_{1}$ has a continuous density, every measurable set is a $\sigma_{1}$-continuous set. This proves the theorem. \(\square\)


\end{document}